% 一维泊松融合教学例：像素0-5，Omega={1,2,3,4}，边界值f*_0=10、f*_5=20，
% 源图像g=(0,0,4,8,4,0)；直接粘贴(10,0,4,8,4,20)，泊松解(10,12,18,24,22,20)。宽约93mm。
\begin{tikzpicture}[x=13mm,y=1.6mm,font=\figurefont\small]
 \fill[BookPaper] (.5,-1) rectangle (4.5,27);
 \node[anchor=north west,text=BookMuted] at (.55,27) {区域$\Omega$};
 \draw[BookMuted,line width=.7pt,->] (-.3,0)--(5.5,0) node[right] {像素};
 \draw[BookMuted,line width=.7pt,->] (-.3,0)--(-.3,28) node[above] {值};
 \foreach \y in {10,20} {
  \draw[BookRule,line width=.5pt] (-.3,\y)--(5.3,\y);
  \node[anchor=east] at (-.32,\y) {$\y$};
 }
 \node[anchor=east] at (-.32,0) {$0$};
 \foreach \x in {0,...,5} {
  \draw[BookMuted,line width=.5pt] (\x,0)--(\x,-.8);
  \node[below] at (\x,-.8) {$\x$};
 }
 \draw[BookBlue,dashed,line width=1pt] plot[mark=*,mark size=1.3pt,
  mark options={solid}] coordinates {(0,10)(1,0)(2,4)(3,8)(4,4)(5,20)};
 \draw[BookTeal,line width=1pt] plot[mark=*,mark size=1.3pt] coordinates {
  (0,10)(1,12)(2,18)(3,24)(4,22)(5,20)};
 \fill[BookMuted] (-.08,9.2) rectangle (.08,10.8);
 \fill[BookMuted] (4.92,19.2) rectangle (5.08,20.8);
 \draw[BookBlue,dashed,line width=1pt] (.2,-8.5)--(.6,-8.5);
 \node[anchor=west] at (.65,-8.5) {直接粘贴};
 \draw[BookTeal,line width=1pt] (2.2,-8.5)--(2.6,-8.5);
 \node[anchor=west] at (2.65,-8.5) {泊松融合};
 \fill[BookMuted] (4.05,-9.5) rectangle (4.25,-7.5);
 \node[anchor=west] at (4.3,-8.5) {边界值};
\end{tikzpicture}
